% Generated from locale/tr/content/incompleteness/theories-computability/first-incompleteness.tex; do not hand-edit.


\olfileid[tr]{inc}{tcp}{inc}

\olsection{$\Th{Q}$'nun Tam, Tutarlı, !!^{axiomatizable}
  Genişletmesi Yoktur}

\begin{thm}
\ollabel{thm:first-incompleteness}
$\Th{Q}$'nun tam, tutarlı, !!{axiomatizable} hiçbir genişletmesi yoktur.
\end{thm}

\begin{proof}
$\Th{Q}$'nun tutarlı, karar verilebilir hiçbir genişletmesi olmadığını zaten
biliyoruz. Oysa $\Th{T}$ tam ve !!{axiomatizable} ise karar verilebilirdir.
\end{proof}

\begin{explain}
Bu teorem, G\"odel'in 1931'de verdiği Birinci Eksiklik Teoremi'nin özgün
ifadesinden pek uzak değildir. Daha çağdaş terminoloji bir yana, temel
farklar şunlardır: G\"odel “tutarlı” yerine “$\omega$-tutarlı” der; ayrıca
biçimsel hesaplanabilirlik kavramı henüz ortaya konmamış olduğu için tam
genellik içinde “!!{axiomatizable}” diyemezdi. (Hesaplanabilirliğin biçimsel
modelleri izleyen on yılda, aralarında G\"odel de bulunan araştırmacılarca,
büyük ölçüde G\"odel olgusuna tabi kuram türlerini niteleyebilmek amacıyla
geliştirilmiştir.)

Teorem, her şeye birden---yani tamlığa, tutarlılığa ve
!!{axiomatizability} özelliğine---sahip olamayacağınızı söyler. Bununla
birlikte bunlardan herhangi birinden vazgeçerseniz öteki ikisine sahip
olabilirsiniz: $\Th{Q}$ tutarlıdır ve hesaplanabilir biçimde
aksiyomlaştırılmıştır, fakat tam değildir; tutarsız kuram tamdır ve
hesaplanabilir biçimde aksiyomlaştırılmıştır (örneğin
$\{ \eq/[0][0] \}$ ile), fakat tutarlı değildir; aritmetiğin doğru tümceleri
kümesi ise tam ve tutarlıdır, fakat hesaplanabilir biçimde
aksiyomlaştırılmamıştır.
\end{explain}
